(1) 对可逆绝热理想气体有 $TV^{\gamma-1}=\mathrm{const}$。压缩 $V_1/V_2=r$，所以 $T_2=T_1r^{\gamma-1}$；膨胀 $V_3/V_4=1/r$，所以 $T_4=T_3/r^{\gamma-1}$，且 $V_3=V_2,V_4=V_1$。
(2) 四顶点温度为 $T_1,\ T_1r^{\gamma-1},\ T_3,\ T_3r^{1-\gamma}$。定容吸热 $Q_{\rm in}=C_V(T_3-T_2)>0$，放热的正大小 $Q_{\rm out}=C_V(T_4-T_1)>0$；因完整循环 $\Delta U=0$，净对外功 $W_{\rm net}=Q_{\rm in}-Q_{\rm out}$。
(3) 利用 $T_4-T_1=(T_3-T_2)/r^{\gamma-1}$，得 $Q_{\rm out}/Q_{\rm in}=r^{1-\gamma}$，故热效率 $\eta=W_{\rm net}/Q_{\rm in}=1-r^{1-\gamma}$。在恒定 $\gamma$ 的空气标准模型中，效率只显式依赖压缩比而不直接取决于 $T_3$，但净功大小仍与吸热温差相关。
(4) $r\to1^+$ 时 $\eta\to0$，压缩比增大使 $\eta$ 单调增加且理论极限趋一；给定 $T_1,T_3$ 要求 $T_3>T_1r^{\gamma-1}$ 才有正输入热量。实际发动机受爆震、摩擦、传热、可变热容和材料限制，不能将无限压缩比视为可行设计。
